\section{Introdução}

Inicialmente, os robôs utilizavam sistemas de controle de alto nível para
evitar colisões; porém, essa abordagem nem sempre é prática quando o robô
precisa reagir rapidamente durante o deslocamento. Por conta disso,      
surgiu, na década de 1980, o controle de baixo nível, que permitiu aos   
robôs realizar o desvio de obstáculos em tempo real~\cite{khatib1986}.   
Esse método possibilitou que o robô interagisse com o ambiente de forma  
mais rápida por meio de sensores acoplados, permitindo atualmente que    
robôs aspiradores mais simples limpem o ambiente de forma autônoma sem   
utilizar um mapa~\cite{corke2023}. Por outro lado, o planejamento de     
caminho baseado em mapa (\textit{path planning}) exige construir ou      
conhecer o mapa do ambiente e saber a posição do robô em relação a ele em
todos os momentos, sendo mais comum em carros autônomos, aeronaves, entre
outros, por suportar tarefas mais complexas~\cite{corke2023}. Já para o  
exemplo do robô aspirador, a abordagem reativa torna-se adequada, pois   
sistemas reativos costumam ser mais rápidos e simples~\cite{corke2023}.  

Mesmo sem mapa global, um robô diferencial precisa manter uma estimativa 
contínua da própria pose com base na velocidade, direção e intervalo de  
deslocamento a partir de uma estimativa anterior~\cite{corke2023}.       
Existem vários métodos para realizar esse cálculo da odometria, sendo um 
deles o método de Runge-Kutta de 2ª~ordem (RK2). Para auxiliar o robô    
durante a navegação, utilizam-se sensores como os \textit{encoders} nas  
rodas (para a odometria) e o LiDAR (para a percepção local dos           
obstáculos); todavia, esses sensores não são perfeitos e acumulam        
erros~\cite{siegwart2004,corke2023}. Nesse contexto, empregam-se         
algoritmos reativos como os Campos Potenciais Artificiais
(APF)~\cite{khatib1986} e o algoritmo Bug2~\cite{lumelsky1987}, que, a   
partir da odometria estimada e das leituras locais, conduzem o robô até a
meta desviando de obstáculos.

Diante disso, o presente trabalho tem como objetivo avaliar
comparativamente os algoritmos clássicos APF e Bug2, utilizando a        
odometria calculada via RK2, em três cenários: o primeiro com obstáculos 
esparsos e um corredor até o ponto final (Cenário~A); o segundo com um   
obstáculo côncavo em formato de ``U'' antes do ponto final (Cenário~B); e
o terceiro com uma barreira posicionada atrás do ponto final (subcenário 
GNRON). Além disso, avalia-se a sensibilidade dos parâmetros do APF e    
valida-se o estimador odométrico RK2 frente ao \textit{ground truth}.

O restante deste artigo está organizado da seguinte forma: a Seção~II  
apresenta a fundamentação teórica e os trabalhos relacionados; a     
Seção~III descreve a metodologia experimental; a Seção~IV detalha o     
desenvolvimento e a implementação; a Seção~V apresenta os resultados
experimentais; a Seção~VI traz a discussão técnica e as limitações; e a 
Seção~VII apresenta as conclusões. 
